\origpage{326}{N. Robertson, P.D. Seymour / Journal of Combinatorial Theory, Series B 90 (2004) 325--385}
\noindent drawn touching the boundary (roughly, we ask that when we do contractions in producing a minor, any such edge remain in contact with the boundary during the contraction process).

Our approach is by a grand induction on the complexity of the surface and the well-quasi-order $\Omega$. Thus, we first assume the result for all simpler surfaces, that is, with fewer handles and crosscaps, even if the surface boundary has more components (``cuffs'') and the well-quasi-order is bigger. Second, for a fixed number of handles and crosscaps, we assume the result for sets of labelled hypergraphs in which the labels of the ``internal'' edges (that is, not on the boundary of $\Sigma$) all come from some proper subideal of our well-quasi-order, even if the labels on the boundary come from some larger well-quasi-order. Third, we proceed by induction on the number of cuffs, and there is a fourth of the same kind. To make this induction work, we find it necessary to divide the boundary of our surface into segments, and then have different restrictions on the labels of edges bordering each segment; and also, some edges drawn on the boundary have to be regarded as fixed.

The method of proof is to apply two theorems of earlier papers concerned with ``tangles'' in hypergraphs. Let $H_1, H_2, \dots$ be a ``bad'' sequence of hypergraphs all drawn in $\Sigma$, with labelling functions $\phi_1, \phi_2, \dots$, such that each pair $(H_i, \phi_i)$ satisfies the restrictions on labels described above. A theorem of [7] implies that if $\mathscr{T}$ is a tangle in $H_i$, and $(H_i, \phi_i)$ is ``sufficiently general'' relative to $\mathscr{T}$, then $(H_i, \phi_i)$ contains $(H_1, \phi_1)$ in the required way, a contradiction. It follows that, for every tangle $\mathscr{T}$ in every $H_i$, there is one of a bounded number of ``structural deficiencies''. But a theorem of [10] says that if relative to every tangle in every $H_i$ we have a suitable kind of decomposition of $H_i$, then again some $(H_j, \phi_j)$ contains some $(H_i, \phi_i)$. Thus, it remains to show that our structural deficiencies can be converted to the right kind of decompositions, and that is the main part of the proof. This is where the grand induction is used---since the pieces into which we propose to decompose our hypergraphs are simpler than the originals (that is, are covered by the inductive hypothesis), we can infer that these pieces satisfy the theorem; and inferring that is a large part of showing that our decompositions are indeed of the ``right kind''.

The paper is organized as follows. We begin in Sections 2 and 3 with definitions, and in Section 4 explain the overall induction. In Sections 5 and 6 we reconcile our containment relation with that of the theorem of [10] (unfortunately, they are not quite the same). Sections 7 and 8 are more definitions, introducing ``tangles'' and ``tie-breakers'', and in Section 9 we state the theorem of [10] that we wish to apply, and begin the application. This is continued in Sections 10--14. Finally, in Section 15 we use a theorem of [7] to complete the proof.

\sectionhead{2. Surfaces}

In this paper, we mean by a \textit{surface} a compact, connected 2-manifold $\Sigma$, with (possibly null) boundary, denoted by $bd(\Sigma)$. An \textit{$O$-arc} in $\Sigma$ is a subset of $\Sigma$ homeomorphic to a circle, and a \textit{line} is a subset of $\Sigma$ homeomorphic to $[0, 1]$. A
